\documentclass[11pt,a4paper]{article}

% ==========================================
% 1. CONFIGURACION GENERAL
% ==========================================

\usepackage[utf8]{inputenc}
\usepackage[T1]{fontenc}
\usepackage[spanish,es-nodecimaldot]{babel}
\usepackage{geometry}
\usepackage{graphicx}
\usepackage{float}
\usepackage{booktabs}
\usepackage{array}
\usepackage{caption}
\usepackage{hyperref}
\usepackage{xcolor}
\usepackage{setspace}

% Margenes y espaciado
\geometry{
    left=2.5cm,
    right=2.5cm,
    top=2.5cm,
    bottom=2.5cm
}

\setstretch{1.15}

% Configuracion de enlaces
\hypersetup{
    colorlinks=true,
    linkcolor=black,
    urlcolor=blue,
    citecolor=black
}

% ==========================================
% 2. PORTADA
% ==========================================

\begin{document}

\begin{titlepage}
    \centering

    {\large Universidad de Buenos Aires\par}
    \vspace{0.3cm}
    {\large Facultad de Ciencias Economicas\par}

    \vspace{1.5cm}

    {\large Tecnicatura Universitaria en
    Gestion y Analisis de Datos en Organizaciones\par}

    \vspace{1.5cm}

    {\Large\bfseries TRABAJO PRACTICO GRUPAL\par}

    \vspace{0.8cm}

    {\LARGE\bfseries
    ¿Es verdad que las canciones globales
    se volvieron mas cortas a lo largo del tiempo?\par}

    \vspace{1cm}

    {\large Informe de analisis de datos\par}

    \vfill

    \begin{flushleft}
        \textbf{Materia:} Laboratorio de Metodos Cuantitativos

        \vspace{0.3cm}

        \textbf{Integrantes:}
        \begin{itemize}
            \item Rocío Ojeda -- OjedaRocio
            \item Valentina Fernandez Zurita -- valentinafernandezu
            \item Hilda Anabel Martinez Uriarte -- hildaanabel
            \item Tania Andrea Campos Aliaga -- tania-1604
        \end{itemize}

        \textbf{Dataset:} track\_data\_final.csv

        \textbf{Repositorio GitHub:}
        \url{https://github.com/OjedaRocio/TP-Grupal-Laboratorio-2C-2026}

        \textbf{Fecha de entrega:} 6 de noviembre de 2026
    \end{flushleft}

\end{titlepage}

% ==========================================
% 3. INFORME ACADEMICO
% ==========================================

\section{Resumen}

El presente trabajo busca analizar la evolución de la duración de las canciones globales a lo largo del tiempo, con el objetivo de determinar si existe una tendencia hacia canciones cada vez más cortas. A partir del análisis de los datos disponibles, se examinarán las variaciones en la duración promedio por año y las diferencias entre distintos períodos, con el fin de identificar patrones temporales y evaluar si los resultados respaldan la hipótesis planteada.

\section{Introduccion}

La duracion de las canciones constituye una caracteristica
relevante para estudiar los cambios en la produccion musical.
La posible reduccion de la duracion a lo largo del tiempo
puede relacionarse con transformaciones en los formatos de
consumo, las estrategias de publicacion y las plataformas
digitales. Sin embargo, estas explicaciones requieren
evidencia adicional y no pueden inferirse unicamente a
partir de la duracion.



La pregunta de investigacion es:
\textit{¿Es verdad que las canciones globales se volvieron
mas cortas a lo largo del tiempo?}



\section{Datos y metodologia}

\subsection{Descripcion del dataset}

El analisis utiliza el archivo
\texttt{track\_data\_final.csv}. Entre las variables de
interes se encuentran la fecha de lanzamiento del album
(\texttt{album\_release\_date}) y la duracion de la pista
(\texttt{track\_duration\_ms}).

La fecha se utiliza para construir una variable temporal,
mientras que la duracion se convierte de milisegundos
a minutos mediante la siguiente expresion:

\[
\text{Duracion en minutos}
=
\frac{\text{Duracion en milisegundos}}{60000}
\]

\subsection{Preparacion de los datos}

Se examinan los tipos de datos, los valores faltantes,
los registros duplicados y los valores atipicos. Las
transformaciones aplicadas se documentan indicando su
justificacion y su efecto sobre la muestra.

\subsection{Estrategia de analisis}

Se calcula la duracion promedio por ano y por periodo.
Ademas, se comparan medidas descriptivas y se representan
graficamente los resultados. La cantidad de observaciones
disponibles en cada periodo se considera al interpretar
las diferencias.

La fecha del album se utiliza como aproximacion temporal
cuando no se dispone de la fecha original de lanzamiento
de cada pista.

\section{Resultados}

\subsection{Evolucion temporal de la duracion}

La Figura \ref{fig:evolucion} presenta la evolucion de la
duracion promedio por ano. El grafico permitira evaluar
si la serie muestra una tendencia descendente, ascendente
o sin una direccion clara.

\begin{figure}[H]
    \centering
    \IfFileExists{evolucion_duracion.png}{
        \includegraphics[width=0.9\textwidth]
        {evolucion_duracion.png}
    }{
        \fbox{
            \parbox[c][4cm][c]{0.85\textwidth}{
                \centering
                Insertar aqui el grafico
                evolucion\_duracion.png
            }
        }
    }
    \caption{Evolucion de la duracion promedio por ano.}
    \label{fig:evolucion}
\end{figure}

La interpretacion de este grafico se completara a partir
de los resultados obtenidos en Python, considerando
posibles fluctuaciones y diferencias en el numero de
observaciones por ano.

\subsection{Distribucion de la duracion}

La Figura \ref{fig:distribucion} muestra la distribucion
de la duracion de las canciones. Su analisis permite
examinar la concentracion de los valores, la dispersion
y la presencia de observaciones extremas.

\begin{figure}[H]
    \centering
    \IfFileExists{distribucion_duracion.png}{
        \includegraphics[width=0.85\textwidth]
        {distribucion_duracion.png}
    }{
        \fbox{
            \parbox[c][4cm][c]{0.8\textwidth}{
                \centering
                Insertar aqui el grafico
                distribucion\_duracion.png
            }
        }
    }
    \caption{Distribucion de la duracion de las canciones.}
    \label{fig:distribucion}
\end{figure}

Los resultados descriptivos permitiran evaluar si el
promedio constituye una medida representativa de la
duracion de las canciones incluidas en la muestra.

\subsection{Cantidad de observaciones por periodo}

La Figura \ref{fig:observaciones} presenta la cantidad
de canciones analizadas en cada periodo. Esta informacion
es importante para contextualizar las comparaciones y
detectar periodos con una representacion limitada.

\begin{figure}[H]
    \centering
    \IfFileExists{observaciones_periodo.png}{
        \includegraphics[width=0.85\textwidth]
        {observaciones_periodo.png}
    }{
        \fbox{
            \parbox[c][4cm][c]{0.8\textwidth}{
                \centering
                Insertar aqui el grafico
                observaciones\_periodo.png
            }
        }
    }
    \caption{Cantidad de canciones por periodo analizado.}
    \label{fig:observaciones}
\end{figure}

\subsection{Resumen estadistico}

La Tabla \ref{tab:resumen} se completara con los
estadisticos calculados en el notebook.

\begin{table}[H]
    \centering
    \caption{Duracion de las canciones por periodo.}
    \label{tab:resumen}
    \begin{tabular}{lrrr}
        \toprule
        Periodo & Cantidad & Promedio (min) & Mediana (min) \\
        \midrule
        Periodo 1 & -- & -- & -- \\
        Periodo 2 & -- & -- & -- \\
        Periodo 3 & -- & -- & -- \\
        \bottomrule
    \end{tabular}
\end{table}

La comparacion entre periodos permitira establecer si
las diferencias observadas son consistentes con una
reduccion de la duracion promedio.

\section{Conclusiones}

La conclusion final debera responder explicitamente
la pregunta de investigacion a partir de los resultados
obtenidos.

Se indicara si los datos muestran evidencia de una
reduccion de la duracion, si la tendencia no resulta
clara o si los resultados son insuficientes para sostener
una conclusion.

Entre las limitaciones a considerar se encuentran la
representatividad de la muestra, la disponibilidad de
fechas validas, la posible desigualdad en la cantidad
de observaciones por periodo y la diferencia entre la
fecha de lanzamiento del album y la de la pista.

Como extension futura, podria analizarse una fuente
que identifique explicitamente las canciones de alcance
global y permita contrastar los resultados con una
muestra mas representativa.

\section{Reflexion metodologica}

Se utilizo inteligencia artificial como herramienta
de asistencia para organizar el informe, proponer
procedimientos de analisis y mejorar la presentacion
de los resultados.

Las decisiones metodologicas, la verificacion de los
calculos y la interpretacion final deben ser revisadas
por los integrantes del grupo mediante la comparacion
con el notebook ejecutado y los datos originales.

La version definitiva documentara los usos concretos
de IA, las instrucciones empleadas y las revisiones
realizadas por el grupo.

\section{Bibliografia}

\begin{itemize}
    \item Dataset: \textit{track\_data\_final.csv}.
    https://www.kaggle.com/datasets/wardabilal/spotify-global-music-dataset-20092025?resource=download


\end{itemize}

\end{document}